% =============================================================================
% paper/en/main.tex — Canonical English manuscript for AJAE submission.
%
% Title:   [PRE-α3 placeholder] Three-Channel Tenant-Driven Land Transition...
%          ★ Title, abstract, §1 Intro to be REWRITTEN under α3 framing
%            (decision gate A of plan 2026-05-18_section3-alpha3-implementation.md)
% Author:  Lee, Dohyeon (Korea University, Dept. of Food and Resource Economics)
% Status:  §3 α3 rewrite 2026-05-18 (this session) — AHM-extension separability test.
%          Abstract / §1 / §6 Discussion still PRE-α3 framing; flagged with TODO.
%
% Compile: cd paper/en && latexmk -xelatex -interaction=nonstopmode main.tex
% Bib:     ../../Bibliography_base.bib (26 entries, 2026-05-18 CoVe-PASS subset for §3)
%
% Plans:   quality_reports/plans/2026-05-15_paper-section3-bootstrap.md (APPROVED, §3 bootstrap)
%          quality_reports/plans/2026-05-15_section3-precision-rewrite.md (PR #4 precision)
%          quality_reports/plans/2026-05-18_section3-alpha3-implementation.md (THIS SESSION)
% ADRs:    quality_reports/decisions/2026-05-18_theoretical-scope.md (ADR-0001)
%          quality_reports/decisions/2026-05-18_outcome-hierarchy.md (ADR-0002)
%          quality_reports/decisions/2026-05-18_section3-structure.md (ADR-0003)
% =============================================================================

\documentclass[12pt, a4paper]{article}
% =============================================================================
% paper/en/preamble_en.tex — XeLaTeX preamble for the English-canonical paper.
%
% Adapted from Preambles/header.tex (Beamer-shared) with article-class
% additions for AJAE submission target (double-spaced, 12pt, 1in margins).
%
% Bibliography style: natbib author-year (plainnat fallback if aer.bst absent).
% =============================================================================

% Engine detection + encoding (pdfTeX fallback)
\usepackage{iftex}
\ifPDFTeX
  \usepackage[utf8]{inputenc}
\fi

% Math + tables + graphics
\usepackage{amsmath, amssymb, amsthm}
\usepackage{booktabs}
\usepackage{graphicx}
\usepackage{caption}
\captionsetup{font=small, labelfont=bf}

% Page geometry (AJAE convention: 1in margins, A4)
\usepackage[a4paper, margin=1in]{geometry}

% Double-spacing for journal submission
\usepackage{setspace}
\doublespacing

% Palette (mirrors Quarto/theme-template.scss; consistent with Preambles/header.tex)
\usepackage{xcolor}
\definecolor{primary-blue}{HTML}{012169}
\definecolor{primary-gold}{HTML}{B9975B}
\definecolor{highlight-yellow}{HTML}{F2A900}
\definecolor{light-bg}{HTML}{E8EDF5}
\definecolor{jet}{HTML}{1A1A1A}
\definecolor{positive}{HTML}{15803D}
\definecolor{negative}{HTML}{B91C1C}
\definecolor{neutral}{HTML}{525252}
\colorlet{accent}{primary-blue}
\colorlet{accent2}{primary-gold}

% Hyperref (loaded after xcolor so link colors resolve)
\usepackage{hyperref}
\hypersetup{
  colorlinks = true,
  linkcolor  = primary-blue,
  urlcolor   = primary-blue,
  citecolor  = primary-blue,
  pdfborder  = {0 0 0}
}

% Citation engine — author-year style.
% Try plainnat (always available in natbib); aer.bst can be substituted at submission.
\usepackage[round, authoryear]{natbib}
\bibliographystyle{plainnat}

% Convenience macros (re-used from Preambles/header.tex)
\newcommand{\muted}[1]{\textcolor{neutral}{#1}}
\newcommand{\key}[1]{\textcolor{primary-gold}{\textbf{#1}}}
\newcommand{\good}[1]{\textcolor{positive}{#1}}
\newcommand{\bad}[1]{\textcolor{negative}{#1}}

% Section-numbering depth
\setcounter{secnumdepth}{3}
\setcounter{tocdepth}{2}

% Theorem-like environments (Prediction = numbered claim by channel)
\theoremstyle{plain}
\newtheorem{prediction}{Prediction}[section]

% Bib-key TENTATIVE warning footer (visual reminder during draft)
\newcommand{\tentativecite}[1]{\textcolor{negative}{\citep{#1}}}

% Increase tolerance for line-breaking to eliminate small overfull hboxes
% caused by long \texttt{} paths, citation keys, and hyphen-resistant math.
\setlength{\emergencystretch}{3em}
\tolerance=1000
\hbadness=2000


% α3 framework: cross-document refs into online_appendix.pdf
\usepackage{xr-hyper}
\externaldocument{online_appendix}

\title{Three-Channel Tenant-Driven Land Transition:\\
       Evidence from Korea's Per-Farm Flat-Rate Direct Payment Scheme}
\author{Lee, Dohyeon\\
        Korea University \\
        Department of Food and Resource Economics}
\date{Draft: \today}

\begin{document}
\maketitle

\begin{abstract}
\noindent
\textit{[★ TODO 2026-05-18: Abstract is PRE-α3 framing — to be rewritten under
α3 separability-test framing per ADR-0002 in a follow-up plan. Placeholder
text retained verbatim from PR \#4 until that rewrite:]} This paper exploits
the 0.5 ha eligibility cutoff in Korea's Public-Interest Direct Payment Scheme
(PIDPS, effective 2020) to identify the causal effects of a per-farm flat-rate
direct payment on smallholder farm behavior. Using a difference-in-differences
regression-discontinuity design on the Farm Household Economic Survey panel
(2018--2022, $N=14{,}474$ farm-years across $3{,}614$ farms), we document a
three-channel tenant-driven land transition operating along bargaining,
composition, and extensive-margin pathways.
\end{abstract}

\noindent\textbf{Keywords:} Agricultural Household Model; regression discontinuity; direct payment; tenancy; rent capitalization; Korea.

\noindent\textbf{JEL:} Q12, Q15, Q18, D13.

\bigskip

% =============================================================================
\section{Introduction}
% =============================================================================
\label{sec:intro}

\textit{[★ TODO 2026-05-18: §1 to be drafted in a follow-up plan (decision
gate A). Under α3 framing the introduction leads with the AHM separability
test, contribution sentence C2 (``we reject AHM separability for Korean small
farms''), and primary outcome $\text{area\_own}$. The PR \#4 stub described a
five-channel framework and three-channel tenant-driven headline; both have
been superseded by ADRs 0001--0003.]}

% =============================================================================
\section{Institutional Context}
% =============================================================================
\label{sec:institutional}

\textit{[STUB — PIDPS-SFFP statute (Act on Public-Interest Direct Payments to Farmers and Fishers, Article 10), 0.5 ha cultivated-area cutoff, Small-Farmer Flat Payment 1{,}200{,}000 KRW/year vs.\ area-proportional 2{,}050{,}000/1{,}970{,}000/1{,}890{,}000 KRW per hectare tier-by-tier. Three pages.]}

% =============================================================================
\section{Theoretical Framework}
% =============================================================================
\label{sec:theory}

% -----------------------------------------------------------------------------
\subsection{Baseline AHM and the Separability Null}
\label{sec:ahm-baseline}
% -----------------------------------------------------------------------------

We test the \citet{SinghSquireStrauss1986_ahm} Agricultural Household Model (AHM) separability null by exploiting Korea's 0.5 ha eligibility cutoff for the per-farm flat-rate Public-Interest Direct Payment Scheme (PIDPS) Small-Farmer Flat Payment (SFFP). Two AHM extensions---\textbf{wealth-biased liquidity} \citep{CarterOlinto2003_liquidity, Kazukauskas2013_disinvestment} and \textbf{implicit-labor supervision} \citep{EswaranKotwal1986_supervision}---yield independent margins through which separability fails. \textit{We reject AHM separability for Korean small farms.}

\paragraph{Setup.} The household maximizes utility over consumption $C$ and leisure $L_\ell$,
\begin{equation}
\max_{C, L_f, L_o, A_{own}, A_{rent}, K} \; U(C, L_\ell),
\label{eq:ahm-objective}
\end{equation}
subject to the integrated farm-household budget constraint,
\begin{equation}
P_c\, C \;=\; P_q\, F(L_f, A, K) \;+\; w\, L_o \;-\; r\, A_{rent} \;-\; p_K\, I(K) \;+\; T,
\label{eq:ahm-budget}
\end{equation}
the time-allocation constraint $L_f + L_o + L_\ell = \bar L$, and the land identity $A = A_{own} + A_{rent}$. The transfer $T = T_{SFFP} = 1{,}200{,}000$ KRW/year enters disposable income if and only if the household is SFFP-eligible (baseline area $A_{2018,i} \le 5{,}000$ m\textsuperscript{2}, equivalently treatment dummy $D_i = 1$).

\paragraph{Separability theorem.} Under complete markets---credit, insurance, labor, land all frictionless and competitively priced---the program in (\ref{eq:ahm-objective})--(\ref{eq:ahm-budget}) admits a recursive factorization \citep{SinghSquireStrauss1986_ahm}. The household first maximizes farm profit
\begin{equation}
\pi^* \;=\; \max_{L_f, A, K} \; \left\{ P_q\, F(L_f, A, K) - w\, L_f - r\, A_{rent} - p_K\, I(K) \right\}
\label{eq:profit-max}
\end{equation}
treating prices as parametric, then maximizes utility
\begin{equation}
\max_{C, L_o} \; U(C, L_\ell) \quad \text{s.t.} \quad P_c\, C \;=\; \pi^* + w(\bar L - L_\ell) + T
\label{eq:utility-max}
\end{equation}
taking $\pi^*$ as exogenous. A lump-sum transfer $T$ enters \emph{only} (\ref{eq:utility-max}); production decisions are unaffected. We therefore obtain the AHM separability null
\begin{equation}
H_0: \quad \frac{\partial A_{own}}{\partial T} \;=\; \frac{\partial L_f}{\partial T} \;=\; \frac{\partial K}{\partial T} \;=\; \frac{\partial L_o}{\partial T} \;=\; 0.
\label{eq:null}
\end{equation}

\paragraph{Sources of separability failure.} \citet{deJanvryFafchampsSadoulet1991_peasant} catalogue four canonical sources of separability failure---missing credit, missing insurance, missing labor markets, missing land markets---each of which replaces a parametric market price in (\ref{eq:profit-max}) with an endogenous shadow price tied to household preferences. The separability-test literature \citep{Benjamin1992_separation, LaFaveThomas2016_markets} has established demographic shocks as instruments for rejecting (\ref{eq:null}); we instead exploit a sharp policy cutoff. Our two extensions in \S\ref{sec:liquidity}--\S\ref{sec:supervision} formalize the credit-imperfection and labor-imperfection margins respectively, with closed-form predictions whose Lagrangian derivations appear in the Online Appendix \S B.1--B.2.

% -----------------------------------------------------------------------------
\subsection{Wealth-Biased Liquidity Relaxation}
\label{sec:liquidity}
% -----------------------------------------------------------------------------

\textbf{Imperfection.} Smallholders face a wealth-dependent credit access function $\rho(W_i)$ with $\rho'(W_i) > 0$: wealthier households obtain credit at lower marginal cost \citep{CarterOlinto2003_liquidity}. This imperfection is binding empirically in the developing- and lower-middle-income agricultural literature \citep{PittKhandker1998_credit, FosterRosenzweig1995_learning}, and---we argue---in the Korean smallholder context, where farms with small baseline asset bases (predominantly tenants) cannot freely borrow against future rental cash flows to finance land purchase.

\textbf{Asset-threshold rule.} For households at the rental--ownership margin, the indifference condition between renting and buying land yields a wealth threshold $W^*(\mathbf{p})$ as a function of prices $\mathbf{p}$ (rent $r$, land price $p_{land}$, interest rate $R$). A household owns land iff
\begin{equation}
W_i + \frac{T_i}{1-\beta} \;\ge\; W^*(\mathbf{p}),
\label{eq:CO-threshold}
\end{equation}
where $T_i = T_{SFFP} \cdot D_i$ is the household-specific subsidy flow and $\beta$ is the discount factor. The subsidy enters as a permanent-income-equivalent shift in disposable wealth, lowering the effective threshold $W^*$ by the amount $T/(1-\beta)$ for SFFP-eligible households.

\textbf{Closed-form predictions} (derivations in Online Appendix \S B.1; \citealp{CarterOlinto2003_liquidity}; \citealp{Kazukauskas2013_disinvestment}):
\begin{equation}
\boxed{\frac{\partial A_{own,i}}{\partial T_{SFFP}} \;=\; -\,\varphi(W_i^*) \cdot \frac{\partial W^*}{\partial T} \;>\; 0,}
\label{eq:CO-1}
\end{equation}
where $\varphi(W^*) > 0$ is the density of household wealth at the threshold. By the wealth-bias slope $\rho'(W) > 0$, the threshold-crossing density is monotone in baseline tenancy: pure tenants ($s_{0,i} = A_{own,2018,i}/A_{2018,i} = 0$) face the largest $\varphi(W^*)$ in the eligibility window, while pure owner-operators ($s_{0,i} = 1$) sit inframarginally above $W^*$ and exhibit zero response by construction. Formally,
\begin{equation}
\boxed{\frac{\partial^2 A_{own,i}}{\partial T_{SFFP}\, \partial (1 - s_{0,i})} \;>\; 0.}
\label{eq:CO-2}
\end{equation}
This monotone-in-baseline-tenancy gradient is the empirical signature uniquely diagnostic of wealth-biased liquidity and cannot arise from a universal income-effect channel (under which all subgroups would respond proportionally).

\textbf{Capital-adjustment sub-prediction.} Conditional on $A_{own}$ pivot, the credit-constrained capital FOC implies an adjustment threshold $\tau$ on physical investment. If $T_{SFFP}/\tau < 1$---i.e., the subsidy is small relative to the discrete capital adjustment scale---no within-household investment crossing occurs and the marginal transfer is absorbed on non-capital margins:
\begin{equation}
\boxed{\frac{\partial \text{op\_cost\_ex\_rent}_i}{\partial T_{SFFP}} \;\le\; 0 \quad \text{if } T_{SFFP}/\tau < 1.}
\label{eq:CO-3}
\end{equation}
This adjustment-threshold logic parallels the $(S,s)$ inaction band of \citet{CaballeroEngel1999_lumpy} cast within an AHM-internal credit constraint; we take the predictions of \citet{CarterOlinto2003_liquidity, Kazukauskas2013_disinvestment} as load-bearing rather than the macro lumpy-investment framework, which is outside our $\alpha$-strict AHM-extension scope.

% -----------------------------------------------------------------------------
\subsection{Implicit Labor with Supervision}
\label{sec:supervision}
% -----------------------------------------------------------------------------

\textbf{Imperfection.} Labor-market completeness fails when hired labor requires costly supervision \citep{EswaranKotwal1986_supervision}. Each hour of hired labor carries a monitoring cost $m > 0$ borne by the farm operator, driving a wedge between the shadow wage of family labor $w_f$ and the market wage $w_m$,
\begin{equation}
w_{f,i} \;=\; w_m + m \cdot \mathbf{1}[\ell_{h,i} > 0],
\label{eq:EK-wedge}
\end{equation}
where $\ell_{h,i}$ is hired labor demand. When family labor is the marginal type ($\ell_{h,i} > 0$), $w_f$ exceeds $w_m$ by the supervision wedge; allocations of family vs.\ hired vs.\ off-farm labor become non-separable from household preferences.

\textbf{Closed-form prediction} (derivation in Online Appendix \S B.2):
\begin{equation}
\boxed{\frac{\partial \text{off\_farm\_income}_i}{\partial T_{SFFP}} \;\ne\; 0 \quad \text{under non-separability,}}
\label{eq:EK-1}
\end{equation}
with sign determined by the comparative statics of the supervision-cost wedge as the budget constraint slackens. The sign is theoretically indeterminate ex ante because the transfer relaxes both the family-time scarcity (favoring more off-farm labor) and the cash constraint that drives supervised hiring (favoring less off-farm labor as own-farm activity expands); empirical estimation must resolve which dominates.

\textbf{Auxiliary status.} Channel~\ref{sec:supervision} is auxiliary to the wealth-biased liquidity channel of \S\ref{sec:liquidity}. Its rejection does not collapse the contribution: the primary $A_{own}$ gradient (\ref{eq:CO-1})--(\ref{eq:CO-2}) survives independently. Rather, a non-zero $\hat\beta_4$ on $\text{off\_farm\_income}$ would constitute a second independent rejection of separability via a labor-imperfection margin, strengthening the joint test.

% -----------------------------------------------------------------------------
\subsection{Unified Predictions and Equilibrium Rent Caveat}
\label{sec:predictions}
% -----------------------------------------------------------------------------

Table~\ref{tab:alpha3-predictions} consolidates the predictions of \S\ref{sec:liquidity}--\S\ref{sec:supervision} and maps each closed-form result to its econometric counterpart in the difference-in-differences regression-discontinuity (DiD-RD) design of \S\ref{sec:identification}.

\begin{table}[ht]
\centering
\caption{Unified predictions of the AHM extensions and their econometric counterparts.}
\label{tab:alpha3-predictions}
\small
\begin{tabular}{p{2.4cm} p{2.6cm} p{3.0cm} p{1.4cm} p{2.4cm}}
\toprule
Channel & Outcome & Reduced form & Econ. $\hat\beta$ & Source \\
\midrule
Liquidity (primary \#1) & $A_{own}$ & $\partial/\partial T_{SFFP} > 0$ & $\beta_1$ & C-O 2003 \\
Liquidity gradient (primary \#1, heterogeneity) & $A_{own} \times (1 - s_0)$ & $\partial^2/\partial T_{SFFP}\, \partial(1-s_0) > 0$ & $\beta_2$ & C-O 2003 \\
Capital adj.\ (primary \#2) & op\_cost\_ex\_rent & $\le 0$ if $T_{SFFP}/\tau < 1$ & $\beta_3$ & K et al.\ 2013 \\
Labor (auxiliary) & off\_farm\_income & $\ne 0$ under non-sep. & $\beta_4$ & E-K 1986 \\
\midrule
\multicolumn{5}{l}{\textit{Ex-theory (B.1) --- aggregate-equilibrium implication, not model-internal:}} \\
Incidence (ex-theory) & unit\_rent\_price & Reduced rental demand $\Rightarrow r \downarrow$ in equilibrium & $\beta_5$ & Floyd 1965; Alston-James 2002 \\
\bottomrule
\end{tabular}
\end{table}

\paragraph{Pre-registered falsification.} We pre-register two falsification predictions consistent with our framework:
\begin{description}
\item[\textbf{F1.}] If the $\hat\beta_2$ interaction is statistically indistinguishable from zero, or if pure owner-operators ($s_0 = 1$) exhibit the same $\hat\beta_1$ as pure tenants ($s_0 = 0$), the wealth-biased liquidity mechanism \citep{CarterOlinto2003_liquidity} is rejected. We would then reframe as a universal income effect under separability and abandon the wealth-bias hypothesis.
\item[\textbf{F2.}] If $\hat\beta_4$ is statistically zero across all bandwidths (T1, T2, T3), the supervision-cost mechanism \citep{EswaranKotwal1986_supervision} is rejected. The wealth-biased liquidity primary channel survives but the contribution narrows to a single margin.
\item[\textbf{F1 + F2 jointly}] would imply that the AHM separability null (\ref{eq:null}) cannot be rejected for the Korean smallholder setting at the 0.5 ha cutoff. The paper then repositions as a precise null estimate---a developed-country complement to the demographic-instrument null estimates in \citet{LaFaveThomas2016_markets}.
\end{description}

\paragraph{Equilibrium rent caveat (B.1 disclosure).} The Korean data exhibit a negative $\hat\beta_5$ on unit\_rent\_price (P3b-2 reports estimates in the range $-130$ to $-48$ KRW/m\textsuperscript{2} across non-pure-owner subgroups at T1) and a rent\_cost pass-through of approximately $-11.1\%$ at T2---a sign opposite to the positive capitalization estimates in \citet{Kirwan2009_rentcap} for the U.S.\ and \citet{BaldoniCiaian2023_eucap} for the EU. These observed rent movements are \emph{consistent with} but \emph{not derived from} our AHM-extension model. The micro mechanism is: rented-area contraction at the household level (an implication of $\partial A_{own}/\partial T_{SFFP} > 0$ from \ref{eq:CO-1} via substitution from rented to owned land) aggregates into a reduced rental-demand schedule. Standard incidence theory \citep{Floyd1965_incidence, AlstonJames2002_incidence} then maps the reduced rental-demand to a negative equilibrium-rent response. We report $\hat\beta_5$ as reduced-form evidence consistent with the model's empirical footprint, not as a primary theoretical contribution; a full equilibrium treatment of the rental market is outside the scope of this paper.

% =============================================================================
\section{Data}
\label{sec:data}
% =============================================================================

\textit{[STUB — Farm Household Economic Survey (FHES) Wave 1, 2018--2022, $N = 14{,}474$ farm-years across $3{,}614$ farms. Variable construction, sample restrictions, descriptive statistics.]}

% =============================================================================
\section{Identification Strategy}
\label{sec:identification}
% =============================================================================

\textit{[STUB — Difference-in-differences regression discontinuity at 0.5 ha. Running variable $rv_{2018} = A_{2018} - 5{,}000$. Three-tier bandwidth T1/T2/T3. Wild cluster bootstrap at household level. Holm step-down. McCrary density test for parallel-trends gate.]}

% =============================================================================
\section{Results}
\label{sec:results}
% =============================================================================

\textit{[★ TODO 2026-05-18: §5 to be reordered under α3 outcome hierarchy
(ADR-0002): §5.1 $A_{own}$ + monotone gradient (primary \#1), §5.2
op\_cost\_ex\_rent (primary \#2), §5.3 off\_farm\_income (auxiliary), §5.4
ex-theory unit\_rent\_price (B.1 evidence). Prior PR \#4 stub described the
three-channel tenant-driven layout; superseded by ADR-0002.]}

% =============================================================================
\section{Robustness}
\label{sec:robustness}
% =============================================================================

\textit{[STUB — Outlier ladder (raw / IHS / Winsorize p1/p99 / Winsorize p0.5/p99.5); Wild cluster bootstrap on 14 headline cells; McCrary density test result ($\Delta = +0.7$pp $< 2$pp threshold); event-study parallel-trends visual; alternative bandwidth specifications.]}

% =============================================================================
\section{Discussion}
\label{sec:discussion}
% =============================================================================

\textit{[★ TODO 2026-05-18: §6 to be drafted under α3 framing. Suggested
content: (i) comparison to demographic-instrument separability tests
\citep{Benjamin1992_separation, LaFaveThomas2016_markets}; (ii) Korea as a
developed-country setting that nonetheless rejects separability; (iii)
mechanism behind the negative $\hat\beta_5$ (per-farm structure preventing
landlord co-capture, B.1 closure step); (iv) potential one-line references
to bargaining/contract alternatives \citep{BanerjeeGertlerGhatak2002_tenancy}
and the macro-investment heterogeneity literature \citep{KhanThomas2008_idiosyncratic}
as out-of-scope but related comparators.]}

% =============================================================================
\section{Conclusion}
\label{sec:conclusion}
% =============================================================================

\textit{[STUB — One paragraph each on: (i) what we learned (rejection of AHM separability via wealth-biased liquidity primary margin); (ii) why the 0.5 ha cutoff matters (policy-induced identification cleaner than demographic instruments); (iii) what remains for future work (Korean-context magnitude of $\hat\beta_4$ supervision channel, dynamic extension of Carter-Olinto threshold rule, paper/ko translation).]}

% =============================================================================

\bibliography{../../Bibliography_base}

\end{document}
